%%%%%%%%%%%%%%%%%%%%%%%%%%%%%%%%%%%%%%%%%%%%%%%%%%%%%%%%%%%%%%%%%%%%%%%%%%%%%%%
%% StuPro B, "Programmierumgebung Offener Antrieb" (POA)
%% Testplan
%% $Id: testplan.tex,v 1.1 2004/03/19 17:57:31 papier Exp $
%%%%%%%%%%%%%%%%%%%%%%%%%%%%%%%%%%%%%%%%%%%%%%%%%%%%%%%%%%%%%%%%%%%%%%%%%%%%%%%
\documentclass[a4paper,titlepage,12pt,ngerman]{scrbook}
\usepackage{../common/header}
\RCSdef $Revision: 1.1 $
\RCSdef $Date: 2004/03/19 17:57:31 $

\newcommand\version{Version 0.1 \xspace}

\begin{document}

%%%%%%%%%%%%%%%%%%%%%%%%%%%%%%%%%%%%%%%%%%%%%%%%%%%%%%%%%%%%%%%%%%%%%%%%%%%%%%%
%% Deckblatt

\begin{titlepage}
\renewcommand{\thefootnote}{\fnsymbol{footnote}}
{\Huge
\raggedright
\textbf{POA} \\
\huge Programmierumgebung offener Antrieb
\rule{\textwidth}{0.75pt}
\par
}
\begin{flushleft}
\normalsize
\version
\vfill

\begin{figure}[htbp]
\begin{center}
\includegraphics[width=15cm]{poa-logo}
\end{center}
\end{figure}

\end{flushleft}
\vfill

{\parindent=0cm
\Huge Testplan
\normalsize
}


\setcounter{footnote}{0}
\end{titlepage}


%%%%%%%%%%%%%%%%%%%%%%%%%%%%%%%%%%%%%%%%%%%%%%%%%%%%%%%%%%%%%%%%%%%%%%%%%%%%%%%
%% Inhaltsverzeichnis

%% \tableofcontents

%%%%%%%%%%%%%%%%%%%%%%%%%%%%%%%%%%%%%%%%%%%%%%%%%%%%%%%%%%%%%%%%%%%%%%%%%%%%%%%
	
\chapter{Die graphische Benutzungsoberfl"ache}

%%%%%%%%%%%%%%%%%%%%%%%%%%%%%%%%%%%%%%%%%%%%%%%%%%%%%%%%%%%%%%%%%%%%%%%%%%%%%%%
%%%%%%%%%%%%%%%%%%%%%%%%%%%%%%%%%%%%%%%%%%%%%%%%%%%%%%%%%%%%%%%%%%%%%%%%%%%%%
\section{Das Hauptfenster}


%%%%%%%%%%%%%%%%%%%%%%%%%%%%%%%%%%%%%%%%%%%%%%%%%%%%%%%%%%%%%%%%%%%%%%%%%%%%%
\subsection{Windows}
\begin{itemize}
  \item Test: Drei Projekte �ffnen. Fenster nebeneinander, �bereinander
  und �berlappend aufrufen
  \item Erwartetes Verhalten: Fenster werden nebeneinander,
  �bereinander oder �berlappend dargestellt
  \item Verhalten in POA 1.0: korrekt
\end{itemize}

%%%%%%%%%%%%%%%%%%%%%%%%%%%%%%%%%%%%%%%%%%%%%%%%%%%%%%%%%%%%%%%%%%%%%%%%%%%%%
\subsection {Help}
\begin{itemize}
  \item Test: Info aufrufen und die Lizenzbedingungen anzeigen
  \item Erwartetes Verhalten: GPL wird angezeigt
  \item Verhalten in POA 1.0: korrekt
\end{itemize}

%%%%%%%%%%%%%%%%%%%%%%%%%%%%%%%%%%%%%%%%%%%%%%%%%%%%%%%%%%%%%%%%%%%%%%%%%%%%%%%
%%%%%%%%%%%%%%%%%%%%%%%%%%%%%%%%%%%%%%%%%%%%%%%%%%%%%%%%%%%%%%%%%%%%%%%%%%%%%
\section{Allgemein}

%%%%%%%%%%%%%%%%%%%%%%%%%%%%%%%%%%%%%%%%%%%%%%%%%%%%%%%%%%%%%%%%%%%%%%%%%%%%%
\subsection{Farbgebung}
\begin{itemize}
  \item Test: 16 Cores erzeugen. Das erste mit 1ns Takt, das zweite
  mit 2ns Takt. das 16te wieder mit 1ns
  \item Erwartetes Verhalten: Alle Cores haben unterschiedliche
  Farben. Das 16te die Farbe des ersten.
  \item Verhalten in POA 1.0: korrekt
\end{itemize}

%%%%%%%%%%%%%%%%%%%%%%%%%%%%%%%%%%%%%%%%%%%%%%%%%%%%%%%%%%%%%%%%%%%%%%%%%%%%%
\subsection{Tooltip}
\begin{itemize}
  \item Test:
  \item Erwartetes Verhalten:
  \item Verhalten in POA 1.0:
\end{itemize}

%%%%%%%%%%%%%%%%%%%%%%%%%%%%%%%%%%%%%%%%%%%%%%%%%%%%%%%%%%%%%%%%%%%%%%%%%%%%%
\subsection{Beschriftung des Layouts}
\begin{itemize}
  \item Test: \begin{enumerate}
    \item Beschriftung ins Projekt einf�gen
    \item Beschriftung verschieben
    \item Beschriftungstext �ndern
    \item Projekt speichern und POA beenden
    \item POA neustarten und Projekt laden
  \end{enumerate}
  \item Erwartetes Verhalten: Die Beschriftung befindet nach dem
  Neustart noch an der richtigen Stelle.
  \item Verhalten in POA 1.0:
\end{itemize}

%%%%%%%%%%%%%%%%%%%%%%%%%%%%%%%%%%%%%%%%%%%%%%%%%%%%%%%%%%%%%%%%%%%%%%%%%%%%%
\subsection{Zoomen}
\begin{itemize}
  \item Test: Projektdarstellung vergr��ern, abspeichern und POA
  beenden. POA starten und Projekt neu laden.
  \item Erwartetes Verhalten: Projekt wird weiterhin vergr��ert dargestellt.
  \item Verhalten in POA 1.0:
\end{itemize}

%%%%%%%%%%%%%%%%%%%%%%%%%%%%%%%%%%%%%%%%%%%%%%%%%%%%%%%%%%%%%%%%%%%%%%%%%%%%%
%%%%%%%%%%%%%%%%%%%%%%%%%%%%%%%%%%%%%%%%%%%%%%%%%%%%%%%%%%%%%%%%%%%%%%%%%%%%%
\section{POA konfigurieren}

%%%%%%%%%%%%%%%%%%%%%%%%%%%%%%%%%%%%%%%%%%%%%%%%%%%%%%%%%%%%%%%%%%%%%%%%%%%%%
\subsection{Raster einstellen}
\begin{itemize}
  \item Test: 
    \begin{enumerate}
      \item Grid einschalten. 
      \item Raster auf 15 stellen. 
      \item Grid  abschalten. 
      \item Raster auf 25 stellen. 
      \item Grid einschalten. 
      \item Raster auf 50 stellen. 
      \item Grid ausschalten. 
      \item POA beenden. 
      \item POA aufrufen. 
      \item Grid einschalten.
    \end{enumerate}
  \item Erwartetes Verhalten:
    \begin{enumerate}
      \item Raster erscheint. 
      \item Rastergr��e = 15  
      \item Raster wird ausgeblendet 
      \item --
      \item Raster erscheint mit Rastergr��e = 25 
      \item Rastergr��e = 50
      \item Raster wird ausgeblendet 
      \item Rastergr��e und Erscheinungsbild werden gesichert
      \item Eisntellungen werden geladen -> Kein Raster ist zu sehen
      \item Grid erscheint mit Rastergr��e 50
    \end{enumerate}
  \item Verhalten in POA 1.0:
    \begin{enumerate}
      \item 
      \item 
      \item 
      \item 
      \item 
      \item 
      \item 
      \item
      \item
      \item
    \end{enumerate}
\end{itemize}

%%%%%%%%%%%%%%%%%%%%%%%%%%%%%%%%%%%%%%%%%%%%%%%%%%%%%%%%%%%%%%%%%%%%%%%%%%%%%
\subsection{Sprache einstellen}

\subsubsection{Wird die Einstellung �bernommen}
\begin{itemize}
  \item Test: Sprache auf deutsch stellen. POA beenden. POA neustarten.
  \item Erwartetes Verhalten: deutsch bleibt nach Neustart eingestellt
  \item Verhalten in POA 1.0: korrekt
\end{itemize}

\subsubsection{Deutsch: Gibt es noch un�bersetzte Stellen in POA?}
\begin{itemize}
  \item Test: Srache auf deutsch stellen. Neu implementierte und ge�nderte Dialoge aufrufen
  \item Erwartetes Verhalten: Texte sind ausschlie�lich in deutsch
  \item Verhalten in POA 1.0:
\end{itemize}

%%%%%%%%%%%%%%%%%%%%%%%%%%%%%%%%%%%%%%%%%%%%%%%%%%%%%%%%%%%%%%%%%%%%%%%%%%%%%
\subsection{Externe Programme einbinden}

\subsubsection{Externen Editor einbinden}
\begin{itemize}
  \item Test: Editor einstellen. CPU anlegen. Quelltext der CPU bearbeiten
  \item Erwartetes Verhalten: Der Editor wird aufgerufen. Der
  Quelltext im Editor ge�ffnet. �nderungen werden �bernommen.
  \item Verhalten in POA 1.0:
\end{itemize}

\subsubsection{Externen Compiler einbinden}
\begin{itemize}
  \item Test: Compiler einstellen. CPU anlegen. Quelltext f�r CPU
  ausw�hlen. Compiler aufrufen.
  \item Erwartetes Verhalten: Der compiler wird aufgerufen udn der
  Quelltext wird compiliert. Compilermeldungen werden auf dem
  Terminalfenster angezeigt.
  \item Verhalten in POA 1.0:
\end{itemize}

%%%%%%%%%%%%%%%%%%%%%%%%%%%%%%%%%%%%%%%%%%%%%%%%%%%%%%%%%%%%%%%%%%%%%%%%%%%%%
\subsection{C Source Template Pfad angeben}
\begin{itemize}
  \item Test: Nicht vorhandenen Pfad eintragen. CPU anlegen. 'Quelltext
  editieren' aufrufen
  \item Erwartetes Verhalten: Fehlermeldung
  \item Verhalten in POA 1.0:
\end{itemize}


%%%%%%%%%%%%%%%%%%%%%%%%%%%%%%%%%%%%%%%%%%%%%%%%%%%%%%%%%%%%%%%%%%%%%%%%%%%%%
\subsection{Download Einstellung}
\begin{itemize}
  \item Test: Offener Antrieb an Com2 bzw. /dev/tty1 anschlie�en. Port
  entsprechend einstellen. Eine CPU compileren und das Binary auf den
  Offenen Antrieb downloaden und mit 'run' ausf�hren
  \item Erwartetes Verhalten: Programm wird ausgef�hrt
  \item Verhalten in POA 1.0:
\end{itemize}


%%%%%%%%%%%%%%%%%%%%%%%%%%%%%%%%%%%%%%%%%%%%%%%%%%%%%%%%%%%%%%%%%%%%%%%%%%%%%
%%%%%%%%%%%%%%%%%%%%%%%%%%%%%%%%%%%%%%%%%%%%%%%%%%%%%%%%%%%%%%%%%%%%%%%%%%%%%
\section{Ein Projekt anlegen, "offnen, speichern}

%%%%%%%%%%%%%%%%%%%%%%%%%%%%%%%%%%%%%%%%%%%%%%%%%%%%%%%%%%%%%%%%%%%%%%%%%%%%%
\subsection{Neues Projekt anlegen}
\begin{itemize}
  \item Test: Projekt anlegen :-)
  \item Erwartetes Verhalten: Projekt wird im gew�hlten Verzeichnis
  angelegt
  \item Verhalten in POA 1.0:
\end{itemize}

%%%%%%%%%%%%%%%%%%%%%%%%%%%%%%%%%%%%%%%%%%%%%%%%%%%%%%%%%%%%%%%%%%%%%%%%%%%%%
\subsection{Projekt speichern}
\begin{itemize}
  \item Test: In zuvor angelegtes Projekt zwei Cores einf�gen
  miteinander verbinden und abspeichern
  \item Erwartetes Verhalten: Projekt wird abgespeichert
  \item Verhalten in POA 1.0:
\end{itemize}

%%%%%%%%%%%%%%%%%%%%%%%%%%%%%%%%%%%%%%%%%%%%%%%%%%%%%%%%%%%%%%%%%%%%%%%%%%%%%
\subsection{Bestehendes Projekt "offnen}
\begin{itemize}
  \item Test: POA schlie�en und neu starten zuvor gespeichertes
  Projekt wieder aufrufen
  \item Erwartetes Verhalten: Projekt wird vollt�ndig geladen
  \item Verhalten in POA 1.0:
\end{itemize}


%%%%%%%%%%%%%%%%%%%%%%%%%%%%%%%%%%%%%%%%%%%%%%%%%%%%%%%%%%%%%%%%%%%%%%%%%%%%%
%%%%%%%%%%%%%%%%%%%%%%%%%%%%%%%%%%%%%%%%%%%%%%%%%%%%%%%%%%%%%%%%%%%%%%%%%%%%%
\section{Layout drucken}
\begin{itemize}
  \item Test: Projekt anlegen mit Beschriftung, Cores, Verk�pfungen  und
  Taktzeiten. Projekt Drucken
  \item Erwartetes Verhalten: Druck entspricht Bildschirmdarstellung
  es werden alle Objekte dargestellt
  \item Verhalten in POA 1.0: korrekt
\end{itemize}


%%%%%%%%%%%%%%%%%%%%%%%%%%%%%%%%%%%%%%%%%%%%%%%%%%%%%%%%%%%%%%%%%%%%%%%%%%%%%
%%%%%%%%%%%%%%%%%%%%%%%%%%%%%%%%%%%%%%%%%%%%%%%%%%%%%%%%%%%%%%%%%%%%%%%%%%%%%
\section{Bl"ocke}

%%%%%%%%%%%%%%%%%%%%%%%%%%%%%%%%%%%%%%%%%%%%%%%%%%%%%%%%%%%%%%%%%%%%%%%%%%%%%
\subsection{Bl"ocke erzeugen}
\begin{itemize}
  \item Test: Core erzeugen, CPU erzeugen 
  \item Erwartetes Verhalten:
  \item Verhalten in POA 1.0:
\end{itemize}

%%%%%%%%%%%%%%%%%%%%%%%%%%%%%%%%%%%%%%%%%%%%%%%%%%%%%%%%%%%%%%%%%%%%%%%%%%%%%
\subsection{Blockkonfiguration allgemein}
\begin{itemize}
  \item Test:
  \item Erwartetes Verhalten:
  \item Verhalten in POA 1.0:
\end{itemize}

\subsubsection{Blockname "andern}
\begin{itemize}
  \item Test:
  \item Erwartetes Verhalten:
  \item Verhalten in POA 1.0:
\end{itemize}

\subsubsection{Block beschreiben}
\begin{itemize}
  \item Test:
  \item Erwartetes Verhalten:
  \item Verhalten in POA 1.0:
\end{itemize}

\subsubsection{Neue Pins erzeugen}
\begin{itemize}
  \item Test:
  \item Erwartetes Verhalten:
  \item Verhalten in POA 1.0:
\end{itemize}

\subsubsection{Pins l"oschen}
\begin{itemize}
  \item Test:
  \item Erwartetes Verhalten:
  \item Verhalten in POA 1.0:
\end{itemize}

\subsubsection{"Anderungen an Pins vornehmen}
\begin{itemize}
  \item Test:
  \item Erwartetes Verhalten:
  \item Verhalten in POA 1.0:
\end{itemize}

\subsubsection{Offset eingeben}
\begin{itemize}
  \item Test:
  \item Erwartetes Verhalten:
  \item Verhalten in POA 1.0:
\end{itemize}

\subsubsection{Takt eingeben}
\begin{itemize}
  \item Test:
  \item Erwartetes Verhalten:
  \item Verhalten in POA 1.0:
\end{itemize}

\subsubsection{Laufzeit eingeben}
\begin{itemize}
  \item Test:
  \item Erwartetes Verhalten:
  \item Verhalten in POA 1.0:
\end{itemize}


%%%%%%%%%%%%%%%%%%%%%%%%%%%%%%%%%%%%%%%%%%%%%%%%%%%%%%%%%%%%%%%%%%%%%%%%%%%%%
\subsection{I/O Block konfigurieren}
Siehe Kapitel {\bf Blockkonfiguration allgemein}.
\begin{itemize}
  \item Test:
  \item Erwartetes Verhalten:
  \item Verhalten in POA 1.0:
\end{itemize}

%%%%%%%%%%%%%%%%%%%%%%%%%%%%%%%%%%%%%%%%%%%%%%%%%%%%%%%%%%%%%%%%%%%%%%%%%%%%%
\subsection{Core konfigurieren}
\begin{itemize}
  \item Test:
  \item Erwartetes Verhalten:
  \item Verhalten in POA 1.0:
\end{itemize}

%%%%%%%%%%%%%%%%%%%%%%%%%%%%%%%%%%%%%%%%%%%%%%%%%%%%%%%%%%%%%%%%%%%%%%%%%%%%%
\subsection{CPU konfigurieren}
\begin{itemize}
  \item Test:
  \item Erwartetes Verhalten:
  \item Verhalten in POA 1.0:
\end{itemize}

\subsubsection{ID eingeben}
\begin{itemize}
  \item Test:
  \item Erwartetes Verhalten:
  \item Verhalten in POA 1.0:
\end{itemize}

\subsubsection{Die CPU-Programmcode schreiben}
\begin{itemize}
  \item Test:
  \item Erwartetes Verhalten:
  \item Verhalten in POA 1.0:
\end{itemize}

\subsubsection{CPU-Programmcode kompilieren}
\begin{itemize}
  \item Test:
  \item Erwartetes Verhalten:
  \item Verhalten in POA 1.0:
\end{itemize}

%%%%%%%%%%%%%%%%%%%%%%%%%%%%%%%%%%%%%%%%%%%%%%%%%%%%%%%%%%%%%%%%%%%%%%%%%%%%%
\subsection{Multiplexer konfigurieren}
\begin{itemize}
  \item Test:
  \item Erwartetes Verhalten:
  \item Verhalten in POA 1.0:
\end{itemize}

%%%%%%%%%%%%%%%%%%%%%%%%%%%%%%%%%%%%%%%%%%%%%%%%%%%%%%%%%%%%%%%%%%%%%%%%%%%%%
\subsection{Blockbibliothek}

\subsubsection{Block in die Blockbibliothek aufnehmen}
\begin{itemize}
  \item Test:
  \item Erwartetes Verhalten:
  \item Verhalten in POA 1.0:
\end{itemize}

\subsubsection{Blocktyp in der Blockbibliothek "andern}
\begin{itemize}
  \item Test:
  \item Erwartetes Verhalten:
  \item Verhalten in POA 1.0:
\end{itemize}

\subsubsection{Block in der Blockbibliothek umbenennen}
\begin{itemize}
  \item Test:
  \item Erwartetes Verhalten:
  \item Verhalten in POA 1.0:
\end{itemize}

\subsubsection{Block aus der Blockbibliothek entfernen}
\begin{itemize}
  \item Test:
  \item Erwartetes Verhalten:
  \item Verhalten in POA 1.0:
\end{itemize}

%%%%%%%%%%%%%%%%%%%%%%%%%%%%%%%%%%%%%%%%%%%%%%%%%%%%%%%%%%%%%%%%%%%%%%%%%%%%%
\subsection{Block in Arbeitsfenster verschieben}
\begin{itemize}
  \item Test:
  \item Erwartetes Verhalten:
  \item Verhalten in POA 1.0:
\end{itemize}

%%%%%%%%%%%%%%%%%%%%%%%%%%%%%%%%%%%%%%%%%%%%%%%%%%%%%%%%%%%%%%%%%%%%%%%%%%%%%
\subsection{Block kopieren und einf"ugen}
\begin{itemize}
  \item Test:
  \item Erwartetes Verhalten:
  \item Verhalten in POA 1.0:
\end{itemize}

%%%%%%%%%%%%%%%%%%%%%%%%%%%%%%%%%%%%%%%%%%%%%%%%%%%%%%%%%%%%%%%%%%%%%%%%%%%%%
\subsection{Block ausschneiden und einf"ugen}
\begin{itemize}
  \item Test:
  \item Erwartetes Verhalten:
  \item Verhalten in POA 1.0:
\end{itemize}

%%%%%%%%%%%%%%%%%%%%%%%%%%%%%%%%%%%%%%%%%%%%%%%%%%%%%%%%%%%%%%%%%%%%%%%%%%%%%
\subsection{Block l"oschen}
\begin{itemize}
  \item Test:
  \item Erwartetes Verhalten:
  \item Verhalten in POA 1.0:
\end{itemize}

%%%%%%%%%%%%%%%%%%%%%%%%%%%%%%%%%%%%%%%%%%%%%%%%%%%%%%%%%%%%%%%%%%%%%%%%%%%%%
%%%%%%%%%%%%%%%%%%%%%%%%%%%%%%%%%%%%%%%%%%%%%%%%%%%%%%%%%%%%%%%%%%%%%%%%%%%%%
\section{Verbindungen}

%%%%%%%%%%%%%%%%%%%%%%%%%%%%%%%%%%%%%%%%%%%%%%%%%%%%%%%%%%%%%%%%%%%%%%%%%%%%%
\subsection{Verbindung routen}

\subsubsection{Bl"ocke miteinander verbinden und l�schen - Default Router}
\begin{itemize}
  \item Test: Einen Input- und einen Output-Block diagonal versetzt erzeugen. InX mit
  OutX verbinden. Eine Verbindung wieder entfernen
  \item Erwartetes Verhalten: Alle Routen �berschneiden sich. Beim
  sp�teren wiederaufrufen (s. u.), fehlt die gel�schte Verbindung
  \item Verhalten in POA 1.0: korrekt
\end{itemize}

\subsubsection{Bl"ocke miteinander verbinden - Smart Router}
\begin{itemize}
  \item Test: Einen neuen Input- und einen neuen Output-Block diagonal versetzt
  erzeugen. InX mit OutX verbinden. F�r jede Verbindung smart Route aufrufen.
  \item Erwartetes Verhalten: Routen �berschneiden sich anschlie�end
  nicht mehr
  \item Verhalten in POA 1.0: korrekt
\end{itemize}

\subsubsection{Verbindung manuell routen und speichern}
\begin{itemize}
  \item Test: Einen neuen Input- und einen neuen Output-Block diagonal versetzt
  erzeugen. InX mit OutX verbinden. Verbindungen manuell
  korrigieren. Projekt speichern und POA beenden. POA neu starten und
  Projekt laden. 
  \item Erwartetes Verhalten: Alle Verbindungen (default, smart, manuell) werden wie zuvor dargestellt
  \item Verhalten in POA 1.0:
\end{itemize}


%%%%%%%%%%%%%%%%%%%%%%%%%%%%%%%%%%%%%%%%%%%%%%%%%%%%%%%%%%%%%%%%%%%%%%%%%%%%%
%%%%%%%%%%%%%%%%%%%%%%%%%%%%%%%%%%%%%%%%%%%%%%%%%%%%%%%%%%%%%%%%%%%%%%%%%%%%%
\section{Scheduling}
\begin{itemize}
  \item Test:
  \item Erwartetes Verhalten:
  \item Verhalten in POA 1.0:
\end{itemize}

%%%%%%%%%%%%%%%%%%%%%%%%%%%%%%%%%%%%%%%%%%%%%%%%%%%%%%%%%%%%%%%%%%%%%%%%%%%%%
%%%%%%%%%%%%%%%%%%%%%%%%%%%%%%%%%%%%%%%%%%%%%%%%%%%%%%%%%%%%%%%%%%%%%%%%%%%%%
\section{Deploy Project}
\begin{itemize}
  \item Test:
  \item Erwartetes Verhalten:
  \item Verhalten in POA 1.0:
\end{itemize}



%%%%%%%%%%%%%%%%%%%%%%%%%%%%%%%%%%%%%%%%%%%%%%%%%%%%%%%%%%%%%%%%%%%%%%%%%%%%%

\end{document}
